\chapter{Dokumentation des KI-Einsatzes und Promptverzeichnis}
\appendixentry{Dokumentation des KI-Einsatzes und Promptverzeichnis}
\label{app:ki-dokumentation}

Gemäß den Vorgaben der Rheinischen Hochschule Köln zur guten wissenschaftlichen Praxis und Transparenz beim Einsatz von generativer künstlicher Intelligenz (KI) wird im Folgenden der genaue Umfang der KI-gestützten Vorarbeiten dargelegt.

\section{Methodische Abgrenzung der Eigenleistung}
Der Einsatz von generativen Sprachmodellen (Google Antigravity / LLM) beschränkte sich strikt auf Hilfs-, Beschleunigungs- und Syntax-Aufgaben:
\begin{itemize}
    \item Erstellung der zustandslosen Demo-Webanwendung (Mock-API und HTML-Dashboard), da die reine Anwendungsentwicklung nicht den wissenschaftlichen Kern der Arbeit darstellt.
    \item Bereitstellung von Python-Messschleifen zur Latenzerfassung.
    \item Syntaxvalidierung von Konfigurationsdateien (GitHub Actions, Kustomize, Terraform).
    \item Typografische Unterstützung bei der Formatierung von LaTeX-Tabellen.
\end{itemize}

Die Konzeption, die Formulierung der Forschungsleitfrage, die architektonischen Entwurfsentscheidungen, die Durchführung der Störungsexperimente auf AWS sowie die kritische Interpretation der Ergebnisse stellen originäre Eigenleistungen des Autors dar.

\section{Promptverzeichnis}
Nachfolgend sind die zehn zentralen Prompts aufgeführt, mit denen die unterstützenden Komponenten generiert und formatiert wurden:

\begin{enumerate}[leftmargin=*]
    \item \promptentry{FastAPI Demo-Backend}{„Schreibe mir ein minimales Python-FastAPI-Backend für einen Kubernetes-Testserver. Ich brauche Standard-Routen für \texttt{/health/live}, \texttt{/health/ready}, \texttt{/api/info} (liefert Hostname, IP und Uptime) sowie einen POST-Endpunkt \texttt{/lab/crash}, der den Prozess hart via \texttt{os.\allowbreak\_exit(1)} beendet.“}
    
    \item \promptentry{Web-Dashboard UI}{„Erstelle mir eine einfache, responsive HTML-Seite mit modernem Dark-Mode-CSS für die FastAPI-App. Die Seite soll per JavaScript-Fetch alle 1,5 Sekunden \texttt{/api/info} abfragen und den antwortenden Pod-Namen sowie die Latenz anzeigen, damit man Pod-Wechsel visuell nachvollziehen kann. Keine externen Frameworks wie React oder Tailwind.“}
    
    \item \promptentry{Gehärtetes Dockerfile}{„Erstelle mir ein schlankes Dockerfile auf Basis von \texttt{python:3.11-slim} für diese FastAPI-App. Achte darauf, dass die App unter einem unprivilegierten Benutzer (\texttt{non-root user}) läuft und ein Standard-HEALTHCHECK definiert ist.“}
    
    \item \promptentry{Lastgenerator-Skript}{„Schreibe mir ein Python-Skript mit \texttt{concurrent.\allowbreak futures.\allowbreak ThreadPoolExecutor}, das kontinuierlich 15 parallele HTTP-GET-Requests pro Sekunde gegen eine URL sendet, HTTP-Statuscodes zählt und am Ende die Perzentile $p_{50}$ und $p_{95}$ der Antwortzeiten als JSON ausgibt.“}
    
    \item \promptentry{GitOps-Drift Automatisierung}{„Erstelle mir ein Python-Skript, das über \texttt{kubectl} den Befehl \texttt{scale\allowbreak{} deployment\allowbreak{} -{}-replicas=0} absetzt und in einer Sekundenschleife den Status \texttt{status.\allowbreak sync.\allowbreak status} der Argo-CD-Application abfragt, um die exakte Zeitspanne bis zum Auftreten von \texttt{OutOfSync} sowie die Zeit bis zum Wiedererreichen von 2 Replikaten zu stoppen.“}
    
    \item \promptentry{GitHub Actions Workflow}{„Wie sieht die YAML-Syntax für eine GitHub-Actions-Pipeline aus, die bei einem Push auf \texttt{main} ein Docker-Image baut und sich mit dem automatischen \texttt{GITHUB\_TOKEN} an der GitHub Container Registry (\texttt{ghcr.io}) anmeldet?“}
    
    \item \promptentry{Kustomize-Overlays}{„Gib mir die Kustomize-Syntax für ein Base-Deployment (1 Replikat) und ein Overlay \texttt{overlays/aws}, das per Patch die Replica-Zahl auf 2 erhöht und eine Traefik-Ingress-Ressource hinzufügt.“}
    
    \item \promptentry{AWS Security Group Parameter}{„Ich erhalte bei Terraform den Fehler: \texttt{InvalidParameterValue:\allowbreak{} Value\allowbreak{} for\allowbreak{} parameter\allowbreak{} GroupDescription\allowbreak{} is\allowbreak{} invalid.\allowbreak{} Character\allowbreak{} sets\allowbreak{} beyond\allowbreak{} ASCII\allowbreak{} are\allowbreak{} not\allowbreak{} supported}. Woran liegt das in meiner \texttt{security-groups.tf}?“}
    
    \item \promptentry{k3s TLS-SAN-Konfiguration}{„Mein lokales \texttt{kubectl} wirft \texttt{x509:\allowbreak{} certificate\allowbreak{} is\allowbreak{} valid\allowbreak{} for\allowbreak{} 10.0.1.88,\allowbreak{} not\allowbreak{} 63.176.27.134}, wenn ich mich mit der Elastic IP verbinde. Wie konfiguriere ich \texttt{tls-san} in k3s nachträglich, damit das Zertifikat die öffentliche IP akzeptiert?“}
    
    \item \promptentry{LaTeX-Tabellengestaltung}{„Ich habe folgende Messergebnisse aus einem Störungsversuch: Zeitpunkte $t_0=0\text{ s}$ (Synced, 2 Pods), $t_2=2{,}6\text{ s}$ (OutOfSync), $t_4=11{,}3\text{ s}$ (Synced, 2 Pods). Formatiere mir diese Daten in sauberem LaTeX-Code unter Verwendung des \texttt{booktabs}-Pakets mit \texttt{\textbackslash toprule}, \texttt{\textbackslash midrule} und \texttt{\textbackslash bottomrule} passend für eine DIN-A4-Projektarbeit.“}
\end{enumerate}